\chapter{Assessment Rubrics}

\section{Weekly Deliverable Rubric}
\begin{longtable}{p{0.18\textwidth}p{0.22\textwidth}p{0.22\textwidth}p{0.22\textwidth}}
\toprule
\textbf{Criterion} & \textbf{Exemplary} & \textbf{Developing} & \textbf{Needs revision}\\
\midrule
\endfirsthead
\toprule
\textbf{Criterion} & \textbf{Exemplary} & \textbf{Developing} & \textbf{Needs revision}\\
\midrule
\endhead
Research focus & Question is specific, feasible, and connected to chemical evidence. & Question is relevant but broad or missing a clear comparison. & Question is vague, untestable, or disconnected from available data.\\
Reproducibility & Scripts run from the project root and outputs are clearly named. & Most work is reproducible but some paths or manual steps remain. & Work cannot be rerun without guessing or manual intervention.\\
Evidence traceability & Claims link to tables, fields, figures, or citations. & Some claims have evidence but links are incomplete. & Claims are mostly unsupported or based on broad impressions.\\
Data quality & Missing values, duplicates, and source coverage are inspected. & Some quality checks are present but incomplete. & Data quality issues are not inspected or documented.\\
Interpretation & Conclusions are cautious, specific, and aligned with evidence. & Interpretation is plausible but overstates some evidence. & Interpretation makes claims the data cannot support.\\
\bottomrule
\end{longtable}

\section{Final Report Rubric}
\begin{longtable}{p{0.18\textwidth}p{0.22\textwidth}p{0.22\textwidth}p{0.22\textwidth}}
\toprule
\textbf{Area} & \textbf{High impact} & \textbf{Acceptable} & \textbf{Revise}\\
\midrule
\endfirsthead
\toprule
\textbf{Area} & \textbf{High impact} & \textbf{Acceptable} & \textbf{Revise}\\
\midrule
\endhead
Abstract & States question, approach, result, limitation, and significance in clear language. & Covers most elements but lacks precision or caveat. & Reads as a general summary without a concrete result.\\
Methods & Includes package version, settings, data source, filtering, validation, and analysis steps. & Includes core workflow but omits some settings or checks. & Too vague to reproduce.\\
Results & Figures and tables directly answer the research question. & Results are relevant but not fully connected to the question. & Results are descriptive without a clear purpose.\\
Figures & Professional, readable, scripted, and captioned with source tables. & Understandable but needs formatting or caption improvements. & Hard to read, manually produced, or not traceable.\\
Discussion & Explains significance, uncertainty, and next steps. & Provides interpretation but limited uncertainty or future direction. & Overclaims or repeats results without interpretation.\\
Reproducibility & Project package is clean, runnable, and documented. & Mostly organized with minor missing files or unclear paths. & Files are disorganized or analysis cannot be rerun.\\
\bottomrule
\end{longtable}

\section{Presentation Rubric}
\begin{tabularx}{\textwidth}{p{0.24\textwidth}Y}
\toprule
\textbf{Criterion} & \textbf{What strong performance looks like}\\
\midrule
Opening & The audience understands the research problem within the first minute.\\
Workflow explanation & The student explains data source, uafR processing, enrichment, and analysis without excessive technical clutter.\\
Main result & The central figure is readable and tied to the research question.\\
Evidence & The student can name the tables, database sources, or citations supporting the result.\\
Limitations & Uncertainty is explicit and scientifically mature.\\
Next steps & The proposed follow-up is realistic and connected to the result.\\
\bottomrule
\end{tabularx}

\section{Review Prompts}
\begin{itemize}
  \item What is the strongest scientific move the student made this week?
  \item What is the most important reproducibility issue to fix?
  \item Which claim needs better evidence?
  \item Which result is ready for a figure, table, or manuscript sentence?
  \item What should the student do before collecting or analyzing more data?
\end{itemize}

\section{Research Log Rubric}
\begin{tabularx}{\textwidth}{p{0.24\textwidth}Y}
\toprule
\textbf{Criterion} & \textbf{Full-credit standard}\\
\midrule
Daily entries & Each workday has an entry with objective, output, evidence inspected, problem, decision, and next step.\\
Decision traceability & Important changes in question, data, filtering, or interpretation are dated and justified.\\
Error documentation & Errors include enough context for another student to understand what failed and how it was fixed.\\
Evidence discipline & Claims in the log refer to source tables, figures, scripts, or citations.\\
Program usefulness & A reviewer can read the log and understand current project status within 10 minutes.\\
\bottomrule
\end{tabularx}

\section{Final Package Rubric}
\begin{longtable}{p{0.20\textwidth}p{0.33\textwidth}p{0.37\textwidth}}
\toprule
\textbf{Component} & \textbf{High-quality evidence} & \textbf{Common deduction}\\
\midrule
\endfirsthead
\toprule
\textbf{Component} & \textbf{High-quality evidence} & \textbf{Common deduction}\\
\midrule
\endhead
README & Explains question, data, scripts, outputs, and rerun order. & Reader cannot tell where to start.\\
Scripts & Numbered, runnable from project root, and write outputs predictably. & Absolute paths, manual steps, or unclear dependencies.\\
Data & Raw data are preserved and processed data are derived by scripts. & Edited raw files or undocumented manual cleaning.\\
uafR outputs & Saved result object, validation summary, source coverage, and trait outputs are present. & Only final figures are saved.\\
Figures & Regenerated by code and captioned with source table and limitation. & Figures are manually edited or overclaim.\\
Report & Methods, results, discussion, limitations, and next steps are complete. & Report lacks reproducibility details or uncertainty.\\
Handoff & Continuation note names what to do next and what not to claim. & Future student would need to reconstruct decisions from scratch.\\
\bottomrule
\end{longtable}
